\scanpage{20}
\mspage{14}

totalement oublier ces ensembles auxiliaires
--- \unsure{j'ai envie|j'aimerais} en tous cas de les oublier.

Mon espoir, c'est que pour une $W$ donnée,
même si on n'arrive pas à trouver
une structure de triple de Quillen
$(W,C,F)$ associée, (p.ex. parce qu'il n'y en
aurait pas), on peut trouver \inserted{à la fois} des
quintuplets \inserted{$(W,C,F,TC,TF)$} qui \unsure{sont|soient} des pré-Quillen à \unsure{g},
et des quintuplets \inserted{(\unsure{d'}\uline{autres} quintuplets \unsure{!|?})} qui \unsure{soient|sont} des pré-Quillen
à droite.

Voici le genre d'énoncé en forme,
pour le moment conjectural voire totalement
heuristique, auquel je songe :

\uline{Théorème-Scholie conjectural} : On suppose
que $M$ est stable par $\varinjlim$ ``quelconques'',
et par $\varprojlim$ finies, et que $W$ est
stable par $\varinjlim$ filtrantes \deleted{\unsure{cela implique}} \inserted{(en plus de la \unsure{rétroversion} faible)}.

\textcircled{1} Alors $\mathrm{Cof}_{W}$ et $\deleted{\unsure{\mathrm{Cof}}}W^{\unsure{\mathrm{univ}}}$ ont les
propriétés de stabilité a) \unsure{et c)|à e)} de la
prop. 2 (p.~4), et aussi la propriété de
stabilité b par facteurs directs, pourvu
que $W$ lui-même l'a.

\textcircled{2} Supposons $M$ \uline{accessible}, de sorte qu'on
dispose du th. 1. \deleted{Pour que} Conditions
équivalentes sur $W$ :
% archive sheet number: 19
